\section{Five open questions (with concrete next steps)}
\label{sec:open-questions}

\subsection*{Q1. Does chromatin-loop topology change which scenario wins?}
\textbf{Basis.} DaMaRiS imposes CTRW sub-diffusion in a homogeneous
2.5\,\textmu m nuclear sphere with a fixed 25\,nm reactive radius.
Real DNA is loop-anchored to cohesin/CTCF-defined TAD boundaries;
two DSBs on the same loop have very different encounter statistics
than two DSBs on distant loops. The paper does not test alternative
polymer/encounter kernels or vary the encounter radius, so
``entwinement'' is a modelling choice rather than a physical
inference.

\textbf{Next steps.} (i)~Re-implement the mean-field surrogate with
a Rouse-polymer inter-DSB distance distribution
($\langle r^2\rangle \sim N^{1/2}$ segment number) and re-rank
Scenarios A--D. (ii)~On free ANL uicgpu A100, run a proper
Geant4-DNA + polymer-chromatin CTRW for one scenario (D) and one
alternative (Rouse-loop), 100 nuclei each; compare residual-DSB
kinetics + mis-rejoin distance distributions. (iii)~Vary encounter
radius 10/25/50\,nm and check whether ranking is preserved.
(iv)~Cross-check against Arnould et al.\ 2021 Hi-C-informed DSB
clustering datasets.

\subsection*{Q2. What happens at high LET where alt-EJ/MMEJ competes?}
\textbf{Basis.} The paper's fit is at 34\,MeV proton, 1.77\,keV/\textmu m
LET, where NHEJ dominates and HR is the main backup. At high LET
(alpha, carbon), complex clustered DSBs shunt to alt-EJ
(Pol\,$\theta$/MMEJ) and slow processive resection pathways not in
the DaMaRiS pathway graph. The four-scenario topology may not be
identifiable at high LET because a missing pathway is a strong
confounder.

\textbf{Next steps.} (i)~Add a fifth pathway state (alt-EJ) to
\texttt{code/damaris\_pathway.py} with published Pol-$\theta$
kinetics (Ceccaldi 2015, Wood 2016). (ii)~Re-run scenario
comparison at 5\,keV/\textmu m and 100\,keV/\textmu m LET, using
Frankenberg-Schwager / Asaithamby 2011 / Reindl 2015 high-LET
foci-vs-time data as reference. (iii)~Test whether Scenario D
remains best or whether an ``entwined + alt-EJ'' hybrid wins at
high LET.

\subsection*{Q3. Can complex chromosome aberrations distinguish the four scenarios?}
\textbf{Basis.} Residual-DSB counts (foci) measure repair
\emph{completion} but not \emph{fidelity}. The paper's entwined
mechanism specifically predicts particular inter-DSB mis-rejoin
statistics --- exactly the observable that would be diagnostic.
mFISH karyotyping and PCC assays yield aberration spectra
(dicentrics, translocations, complex exchanges) that the four
scenarios should predict differently. Not attempted in the paper.

\textbf{Next steps.} (i)~Modify the Gillespie simulator to record
which DSB-end joins which. (ii)~Compare against published mFISH
aberration spectra (Loucas \& Cornforth 2001, Anderson 2002,
Suzuki 2006 --- all publicly tabulated). (iii)~Rank scenarios by
aberration-spectrum log-likelihood, independent of foci-kinetics
fit. (iv)~Quantify the aberration-frequency precision needed to
distinguish D from A.

\subsection*{Q4. Is the conclusion cell-cycle-phase-robust?}
\textbf{Basis.} Beucher 2009 chose G2 to see NHEJ and HR
simultaneously. In G1 (HR unavailable) the four scenarios collapse
to NHEJ-only differences; in S-phase, replication-associated DSBs
use single-ended HR with different kinetics. The claim that
``entwinement is necessary'' may not generalise across the cell
cycle.

\textbf{Next steps.} (i)~Extract cell-cycle-phase-resolved
foci-kinetics data (Riballo 2004, Beucher 2009 G1/S/G2, Kuhne 2004).
(ii)~Re-fit A/B/C/D per phase, allowing HR-branch rates to be
phase-dependent. (iii)~Ask whether Scenario D remains uniquely
best in G1 where entwinement should be least active --- a strong
test of the mechanistic claim.

\subsection*{Q5. Does the model translate to NHEJ-defective cancer cell lines?}
\textbf{Basis.} The paper validates against normal fibroblasts and
MEFs with clean single-gene knockouts. Cancer cells accumulate
multiple partial DDR defects (ATM-low + BRCA1-hypomorphic +
Ku-overexpressing) creating very different pathway-competition
landscapes. A model fit to normal cells may fail to predict repair
phenotypes of BRCA-mutated tumours under PARPi or DNA-PKi --- a
translational gap.

\textbf{Next steps.} (i)~Curate a small (5--10) cancer cell line
panel with published DDR-mutation status and foci-vs-time data
(DepMap + CCLE + individual papers). (ii)~Re-fit Scenario D with
each DDR gene's activity as a per-line rate-constant multiplier.
(iii)~Test whether D uniquely predicts sensitivity to a DDR-inhibitor
panel (PARPi, ATRi, DNA-PKi) vs.\ A/B/C. If D outperforms, that is
strong translational validation; if not, that is a real limit on
the paper's claim of mechanistic generality.
